% ============================================================================
% Stage-1 On-Policy SFT Algorithm Note
% ----------------------------------------------------------------------------
% Input fragment for main.tex. Keep this file free of preamble and document
% environment commands. Shared packages and macros live in config.sty.
% ============================================================================

\section{Two-Stage On-Policy SFT to WDL-SFT Algorithm Note}
\label{sec:opsft-note}
\paragraph{Overview.}
This note formalizes the two-stage experiment used in the current beta-search
and fast-validation cycle. Stage 1 trains a single Qwen3-4B-Base policy with
on-policy supervised fine-tuning (On-Policy SFT): for each prompt, the current
policy samples multiple responses; a verifier assigns binary correctness
labels; correct responses are optimized with ordinary SFT, while incorrect
responses optionally contribute a reverse-SFT term controlled by a scalar
hyperparameter $\beta$. Stage 2 initializes a two-submodel joint policy from
the selected Stage-1 checkpoints, rolls out with Model~2 only, and trains both
submodels through a fused-logit WDL-SFT loss. The purpose of this experiment is
to test whether Stage-1 beta selection and Stage-2 fused-loss gradient
amplification can improve the selected SFT baseline under controlled
conditions.


\section{Problem Setup}
\label{sec:opsft-setup}

Let $x\in\mathcal{X}$ be a math prompt sampled from the training distribution
$\Dtrain$, and let $\pi_{\theta}$ denote the trainable language-model policy.
At training iteration $s$, we sample $N$ responses from the current rollout
policy:
\begin{equation}
    y_i=(y_{i,1},\ldots,y_{i,T_i})
    \sim \pi_{\theta_s}(\cdot\mid x),
    \qquad i=1,\ldots,N .
\end{equation}
The experiment uses $N=8$ rollouts per prompt. A rule-based verifier provides a
binary reward label
\begin{equation}
    r_i = R(x,y_i)\in\{+1,-1\},
\end{equation}
where $+1$ means the response is judged correct and $-1$ means it is judged
incorrect. The responses are partitioned into
\begin{equation}
    \Cset(x)=\{i:r_i=+1\},
    \qquad
    \Iset(x)=\{i:r_i=-1\}.
\end{equation}

For a response $y_i$, define its masked sequence log-likelihood under the
current actor as
\begin{equation}
    \ell_{\theta}(x,y_i)
    =
    \sum_{t=1}^{T_{\max}}
    m_{i,t}
    \log \pi_{\theta}(y_{i,t}\mid x,y_{i,<t}),
    \label{eq:opsft-seq-logprob}
\end{equation}
where $m_{i,t}\in\{0,1\}$ masks out padding tokens. In the implementation this
quantity is computed from \texttt{log\_prob} and \texttt{response\_mask}.

\paragraph{Stage-1 experimental scope.}
The first half of this document describes the Stage-1 run family:
\begin{itemize}[leftmargin=1.5em]
    \item single-model training from Qwen3-4B-Base;
    \item \texttt{loss\_mode=wdl\_sft};
    \item \texttt{joint\_training=False};
    \item \texttt{rollout\_is=null}, \texttt{rollout\_rs=null};
    \item KL in reward and KL loss disabled;
    \item beta grid $\beta\in\{0.0,0.1,\ldots,1.0\}$, with
    $\beta\in\{0.0,0.1,0.2,0.3\}$ completed in the current reported segment.
\end{itemize}
The second half describes the Stage-2 fast validation for the selected
Stage-1 checkpoints. IS-corrected WDL-SFT and group-advantage losses remain
separate methods and are intentionally not part of this note.

\section{Loss Function}
\label{sec:opsft-loss}

The On-Policy SFT loss has two components: a positive SFT term on correct
rollouts and an optional reverse-SFT term on incorrect rollouts.

\subsection{Positive SFT on Correct Rollouts}

When at least one sampled response is correct, the positive term is
\begin{equation}
    \Lpos(x;\theta)
    =
    -\frac{1}{|\Cset(x)|}
    \sum_{i\in\Cset(x)}
    \ell_{\theta}(x,y_i).
    \label{eq:opsft-lpos}
\end{equation}
Minimizing $\Lpos$ increases the likelihood of verifier-correct on-policy
responses. This is ordinary supervised fine-tuning, but the supervision is
created online by the current policy rather than taken from a fixed human
demonstration dataset.

\subsection{Reverse SFT on Incorrect Rollouts}

For incorrect rollouts, define
\begin{equation}
    \Lneg(x;\theta)
    =
    \frac{1}{|\Iset(x)|}
    \sum_{j\in\Iset(x)}
    \ell_{\theta}(x,y_j),
    \qquad |\Iset(x)|>0.
    \label{eq:opsft-lneg}
\end{equation}
This term has the opposite sign from ordinary SFT. Since log probabilities are
typically negative, minimizing $\Lneg$ decreases the likelihood of incorrect
responses. The scalar $\beta\ge 0$ controls how strongly this reverse update is
applied.

\subsection{Total Objective and Boundary Cases}

The per-prompt objective is
\begin{equation}
    \Ltotal(x;\theta)
    =
    \begin{cases}
    \Lpos(x;\theta)+\beta\,\Lneg(x;\theta),
        & |\Cset(x)|>0,\ |\Iset(x)|>0,\\[0.4em]
    \Lpos(x;\theta),
        & |\Cset(x)|>0,\ |\Iset(x)|=0,\\[0.4em]
    0,
        & |\Cset(x)|=0.
    \end{cases}
    \label{eq:opsft-total-loss}
\end{equation}
The batch loss is the same construction applied to the response batch consumed
by the actor update: the implementation first converts verifier scores into raw
reward labels, writes them into the \texttt{advantages} tensor, and then the
loss reads the label of each response from \texttt{advantages[:,0]}. Thus the
\texttt{advantages} tensor in this loss mode is not a GRPO advantage; it is a
carrier for binary reward labels.

The current code normalizes the positive and negative terms by the number of
responses in their respective sets. This keeps the scale of each term less
sensitive to the correct/incorrect ratio inside a mini-batch.

The training objective is the on-policy empirical risk
\begin{equation}
    \min_{\theta}\;
    \E_{x\sim\Dtrain}
    \E_{y_1,\ldots,y_N\sim \pi_{\theta_{\mathrm{old}}}(\cdot\mid x)}
    \left[
        \Ltotal(x;\theta)
    \right],
    \label{eq:opsft-population-objective}
\end{equation}
where the sampled responses and verifier labels are treated as fixed targets
during the actor update. The loss is therefore SFT-like in the inner update,
but on-policy in how the targets are generated.

\section{Algorithm}
\label{sec:opsft-algorithm}

\begin{algorithm}[t]
\caption{Stage-1 On-Policy SFT with reverse-SFT beta}
\label{alg:opsft-on-policy-sft}
\begin{algorithmic}[1]
\Require initial policy $\pi_{\theta_0}$; verifier $R$; beta $\beta$; rollouts per prompt $N$; training prompts $\Dtrain$
\For{training step $s=0,1,\ldots,S-1$}
    \State Sample a batch of prompts $x\sim\Dtrain$.
    \ForAll{prompts $x$ in the batch}
        \State Generate $N$ responses $y_1,\ldots,y_N\sim \pi_{\theta_s}(\cdot\mid x)$.
        \State Score each response with the verifier: $r_i=R(x,y_i)\in\{+1,-1\}$.
        \State Form $\Cset(x)=\{i:r_i=+1\}$ and $\Iset(x)=\{i:r_i=-1\}$.
    \EndFor
    \State Recompute token log-probabilities $\log \pi_{\theta_s}(y_{i,t}\mid x,y_{i,<t})$ with the actor.
    \State Compute $\Ltotal$ from Eq.~\eqref{eq:opsft-total-loss}.
    \State Update $\theta$ by backpropagating $\nabla_{\theta}\Ltotal$.
    \State Every 5 steps, run validation and save latest/best checkpoints.
\EndFor
\end{algorithmic}
\end{algorithm}

\section{Flow Diagram}
\label{sec:opsft-flow}

\begin{figure}[t]
\centering
\begin{tikzpicture}[
    node distance=9mm and 12mm,
    box/.style={
        draw,
        rounded corners=2pt,
        align=center,
        minimum width=30mm,
        minimum height=9mm,
        font=\small
    },
    wide/.style={
        draw,
        rounded corners=2pt,
        align=center,
        minimum width=45mm,
        minimum height=10mm,
        font=\small
    },
    decision/.style={
        draw,
        diamond,
        aspect=2.1,
        align=center,
        inner sep=1pt,
        font=\small
    },
    arrow/.style={-{Latex[length=2.4mm]}, thick}
]
    \node[box] (prompt) {Training prompt\\$x\sim\Dtrain$};
    \node[box, right=of prompt] (rollout) {On-policy rollout\\$y_1,\ldots,y_N\sim\pi_\theta$};
    \node[box, right=of rollout] (verify) {Verifier\\$r_i\in\{+1,-1\}$};

    \node[decision, below=of verify] (split) {Partition\\responses};
    \node[box, below left=13mm and 10mm of split] (correct) {Correct set\\$\Cset$};
    \node[box, below right=13mm and 10mm of split] (incorrect) {Incorrect set\\$\Iset$};

    \node[wide, below=of correct] (pos) {Positive SFT\\$\Lpos=-|\Cset|^{-1}\sum_{i\in\Cset}\ell_\theta(x,y_i)$};
    \node[wide, below=of incorrect] (neg) {Reverse SFT\\$\Lneg=|\Iset|^{-1}\sum_{j\in\Iset}\ell_\theta(x,y_j)$};

    \node[wide, below right=16mm and -21mm of pos] (loss) {Total loss\\$\Ltotal=\Lpos+\beta\Lneg$\\skip if $\Cset=\emptyset$};
    \node[box, below=of loss] (update) {Actor update\\backprop through $\log\pi_\theta$};
    \node[box, below=of update] (eval) {Dense validation\\latest + best checkpoint};

    \draw[arrow] (prompt) -- (rollout);
    \draw[arrow] (rollout) -- (verify);
    \draw[arrow] (verify) -- (split);
    \draw[arrow] (split) -- (correct);
    \draw[arrow] (split) -- (incorrect);
    \draw[arrow] (correct) -- (pos);
    \draw[arrow] (incorrect) -- (neg);
    \draw[arrow] (pos) -- (loss);
    \draw[arrow] (neg) -- (loss);
    \draw[arrow] (loss) -- (update);
    \draw[arrow] (update) -- (eval);
    \draw[arrow] (eval.west) .. controls +(-35mm,0mm) and +(0mm,-18mm) .. (prompt.south);
\end{tikzpicture}
\caption{Current Stage-1 On-Policy SFT training loop. The supervision is generated online: the current policy samples responses, a verifier labels them, and the loss increases correct-response likelihood while optionally decreasing incorrect-response likelihood through $\beta$.}
\label{fig:opsft-flow}
\end{figure}

\section{Interpretation of the Loss}
\label{sec:opsft-interpretation}

\paragraph{Why this is on-policy.}
The responses used for the SFT target are sampled from the current policy
during training. This differs from ordinary offline SFT, where the target
sequence is fixed before training. Here, the model repeatedly proposes its own
candidate solutions and then learns from the subset judged correct by the
verifier.

\paragraph{What the positive term does.}
$\Lpos$ is selected SFT: among the model's own responses, it treats verified
correct responses as demonstrations. If the policy discovers a correct solution
style, the positive term reinforces that behavior.

\paragraph{What the negative term does.}
$\Lneg$ is reverse SFT. It asks the model to reduce probability mass on
verified incorrect responses. This is not a policy-gradient advantage and not
a KL penalty. It is a direct likelihood penalty on complete sampled responses.
The hyperparameter $\beta$ is therefore important: too small may ignore useful
negative evidence, while too large may punish broad regions of the policy and
destabilize useful reasoning formats.

\paragraph{Why all-incorrect prompts are skipped.}
When $\Cset=\emptyset$, the current implementation returns zero loss. This
choice avoids training solely from negative samples when the model has not
found any verified correct answer for the prompt. It makes the algorithm
conservative: every update must be anchored by at least one positive response.

\paragraph{What this loss is not.}
The current Stage-1 loss does not use old/current probability ratios,
importance weights, rollout correction, GRPO advantages, or KL regularization.
These omissions are deliberate for the beta-search phase: the experiment first
isolates the behavior of selected on-policy SFT and the reverse-SFT coefficient
$\beta$.

\section{Beta Search Design}
\label{sec:opsft-beta-search}

The current grid searches only $\beta$ while holding the rest of the training
configuration fixed: same initialization, data seed, learning rate, rollout
count, validation cadence, and checkpoint-selection metric. Each run trains for
150 steps with validation every 5 steps. The within-run best checkpoint is
selected by
\begin{equation}
    \texttt{val-core/HuggingFaceH4/MATH-500/acc/mean@3}.
\end{equation}

\begin{table}[t]
\centering
\small
\setlength{\tabcolsep}{3.3pt}
\renewcommand{\arraystretch}{1.15}
\caption{Stage-1 beta grid. Columns are beta values. Rows are online validation outcomes. Values are percentages from the local metrics JSONL files and the best-checkpoint metadata. The table reports the completed grid segment $\beta\in\{0.0,0.1,0.2,0.3\}$. The MATH-500 checkpoint row is selected by MATH-500 mean@3; rows labeled ``Best'' report each metric's own peak over validation steps.}
\label{tab:opsft-beta-grid}
\resizebox{\textwidth}{!}{%
\begin{tabular}{lcccc}
\toprule
Result / $\beta$ & 0.0 & 0.1 & 0.2 & 0.3 \\
\midrule
Status
    & done & done & done & done \\
Latest step
    & 150 & 150 & 150 & 150 \\
Best checkpoint step
    & 85 & 150 & 65 & 50 \\
Best MATH-500 mean@3
    & 73.25 & \textbf{75.74} & 73.39 & 72.92 \\
Best MATH-500 best@3
    & 80.53 & \textbf{81.66} & 80.20 & 80.42 \\
Final MATH-500 mean@3
    & 70.70 & \textbf{75.74} & 68.28 & 68.55 \\
Best AIME-2025 mean@3
    & 14.10 & 15.38 & 17.95 & \textbf{19.23} \\
Best AIME-2025 best@3
    & 19.25 & 22.10 & 24.24 & \textbf{25.87} \\
Health note
    & stable & stable & MATH drift & MATH drift \\
\bottomrule
\end{tabular}
}
\end{table}

\paragraph{Current reading.}
The completed grid segment suggests a tradeoff. $\beta=0.1$ is the strongest
MATH-500 setting and peaks at the final checkpoint, while larger beta values
increase the best AIME-2025 metrics but drift downward on MATH-500 after their
early best checkpoints. This is exactly the behavior the grid is meant to
measure: reverse-SFT may help some harder generalization metrics, but excessive
negative pressure can reduce the stable MATH trajectory.

\section{Stage-2 Model2-Rollout WDL-SFT}
\label{sec:stage2-wdl-sft}

Stage 2 converts the best Stage-1 single-model checkpoints into the strong
submodel of a two-submodel joint policy. Let $\pi_{\psi}$ denote the Stage-1
policy selected by the beta search. We instantiate a joint model with two
trainable submodels:
\begin{equation}
    \theta=(\theta_1,\theta_2),
    \qquad
    \pi_{\theta_1}^{(1)} \leftarrow \pi_{\mathrm{base}},
    \qquad
    \pi_{\theta_2}^{(2)} \leftarrow \pi_{\psi}.
    \label{eq:stage2-init}
\end{equation}
Both $\theta_1$ and $\theta_2$ are updated during Stage 2. The key distinction
from a fully fused rollout method is that sampling is performed only by
Model~2:
\begin{equation}
    y_i\sim \pi_{\theta_{2,s}}^{(2)}(\cdot\mid x),
    \qquad i=1,\ldots,N.
    \label{eq:stage2-model2-rollout}
\end{equation}
Thus the behavior policy used to generate training targets is the current
Model2 policy. The loss, however, is evaluated on a fused policy distribution
constructed from both submodels' logits.

\subsection{Fused Training Policy}
\label{sec:stage2-fused-policy}

Let $z^{(1)}_{\theta_1,t}(v)$ and $z^{(2)}_{\theta_2,t}(v)$ be the logits of
Model~1 and Model~2 for token $v$ at position $t$ conditioned on
$(x,y_{<t})$. For a fixed mixing coefficient $\lambda\in[0,1]$, the fused
training distribution is
\begin{equation}
    \pi_{\theta}^{\mathrm{mix}}(v\mid x,y_{<t})
    =
    \mathrm{softmax}\!\left(
        (1-\lambda)z^{(1)}_{\theta_1,t}(v)
        +
        \lambda z^{(2)}_{\theta_2,t}(v)
    \right).
    \label{eq:stage2-fused-policy}
\end{equation}
The sequence log-likelihood used by the Stage-2 loss is therefore
\begin{equation}
    \ell_{\theta}^{\mathrm{mix}}(x,y_i)
    =
    \sum_{t=1}^{T_{\max}}
    m_{i,t}
    \log
    \pi_{\theta}^{\mathrm{mix}}(y_{i,t}\mid x,y_{i,<t}).
    \label{eq:stage2-fused-loglik}
\end{equation}
Since $\pi_{\theta}^{\mathrm{mix}}$ depends on both submodels, the gradient
backpropagates through both $\theta_1$ and $\theta_2$. In this design, Model~1
is not frozen. It is allowed to move so that the fused-logit objective acts as
a gradient amplifier for Model~2, while evaluation still targets Model~2
rather than the fused policy.

\subsection{Stage-2 Loss}
\label{sec:stage2-loss}

The verifier and response partition are identical to Stage 1:
\begin{equation}
    r_i=R(x,y_i)\in\{+1,-1\},
    \qquad
    \Cset(x)=\{i:r_i=+1\},
    \qquad
    \Iset(x)=\{i:r_i=-1\}.
\end{equation}
The only loss-level change is the replacement of the single-model
log-likelihood $\ell_{\theta}$ by the fused log-likelihood
$\ell_{\theta}^{\mathrm{mix}}$. The positive and negative Stage-2 terms are
\begin{align}
    \Lpos^{(2)}(x;\theta_1,\theta_2)
    &=
    -\frac{1}{|\Cset(x)|}
    \sum_{i\in\Cset(x)}
    \ell_{\theta}^{\mathrm{mix}}(x,y_i),
    \label{eq:stage2-lpos}
    \\
    \Lneg^{(2)}(x;\theta_1,\theta_2)
    &=
    \frac{1}{|\Iset(x)|}
    \sum_{j\in\Iset(x)}
    \ell_{\theta}^{\mathrm{mix}}(x,y_j).
    \label{eq:stage2-lneg}
\end{align}
The per-prompt Stage-2 objective is
\begin{equation}
    \Ltotal^{(2)}(x;\theta_1,\theta_2)
    =
    \begin{cases}
    \Lpos^{(2)}(x;\theta_1,\theta_2)
    +\beta_2\,\Lneg^{(2)}(x;\theta_1,\theta_2),
        & |\Cset(x)|>0,\ |\Iset(x)|>0,\\[0.4em]
    \Lpos^{(2)}(x;\theta_1,\theta_2),
        & |\Cset(x)|>0,\ |\Iset(x)|=0,\\[0.4em]
    0,
        & |\Cset(x)|=0.
    \end{cases}
    \label{eq:stage2-total-loss}
\end{equation}
The population objective is
\begin{equation}
    \min_{\theta_1,\theta_2}
    \E_{x\sim\Dtrain^{(2)}}
    \E_{y_1,\ldots,y_N\sim \pi_{\theta_{2,\mathrm{old}}}^{(2)}(\cdot\mid x)}
    \left[
        \Ltotal^{(2)}(x;\theta_1,\theta_2)
    \right].
    \label{eq:stage2-objective}
\end{equation}
Here $\Dtrain^{(2)}$ is a non-overlapping shard that starts after the prompts
consumed by the Stage-1 runs under the same data seed. In the current fast
validation, $N=8$, the prompt batch size is 64, and the actor mini-batch size is
512 responses. Consequently, each Model2 rollout batch is consumed by a single
optimizer update with one PPO epoch. This does not remove all sources of
off-policy mismatch---for example, FSDP/vLLM numerical differences still exist
---but it avoids repeatedly reusing the same rollout batch across multiple
mini-batch optimizer updates.

\begin{algorithm}[t]
\caption{Stage-2 Model2-Rollout Fused-Loss WDL-SFT}
\label{alg:stage2-model2-rollout}
\begin{algorithmic}[1]
\Require Stage-1 selected Model2 checkpoint $\pi_{\psi}$; base model $\pi_{\mathrm{base}}$; verifier $R$; mixing weight $\lambda$; Stage-2 beta $\beta_2$; non-overlap prompts $\Dtrain^{(2)}$
\State Initialize $\theta_1\leftarrow \pi_{\mathrm{base}}$ and $\theta_2\leftarrow \pi_{\psi}$.
\For{training step $s=0,1,\ldots,S_2-1$}
    \State Sample a prompt batch $x\sim\Dtrain^{(2)}$.
    \ForAll{prompts $x$ in the batch}
        \State Generate $N$ responses from Model~2 only:
        $y_1,\ldots,y_N\sim\pi_{\theta_{2,s}}^{(2)}(\cdot\mid x)$.
        \State Score each response with $R$ and form $\Cset(x)$ and $\Iset(x)$.
    \EndFor
    \State Run a fused forward pass with both submodels and compute $\pi_{\theta}^{\mathrm{mix}}$ from Eq.~\eqref{eq:stage2-fused-policy}.
    \State Compute $\Ltotal^{(2)}$ from Eq.~\eqref{eq:stage2-total-loss}.
    \State Update both $\theta_1$ and $\theta_2$ by backpropagating $\nabla_{\theta_1,\theta_2}\Ltotal^{(2)}$.
    \State Validate Model~2 and retain latest/best checkpoints.
\EndFor
\end{algorithmic}
\end{algorithm}

\begin{figure}[t]
\centering
\begin{tikzpicture}[
    node distance=8mm and 9mm,
    box/.style={
        draw,
        rounded corners=2pt,
        align=center,
        minimum width=28mm,
        minimum height=8mm,
        font=\small
    },
    wide/.style={
        draw,
        rounded corners=2pt,
        align=center,
        minimum width=43mm,
        minimum height=9mm,
        font=\small
    },
    arrow/.style={-{Latex[length=2.2mm]}, thick}
]
    \node[box] (base) {Qwen3-4B-Base\\initial policy};
    \node[wide, right=of base] (stage1) {Stage 1\\single-model on-policy SFT\\select best checkpoint};
    \node[box, right=of stage1] (model2) {Stage-1 best\\Model~2 init};

    \node[box, below=15mm of base] (m1) {Model~1\\base init\\trainable};
    \node[box, below=15mm of model2] (m2) {Model~2\\Stage-1 init\\trainable};
    \node[wide, below=of stage1] (rollout) {Model2-policy rollout\\$y_i\sim\pi_{\theta_2}^{(2)}(\cdot\mid x)$};
    \node[box, below=of rollout] (verify) {Verifier\\$\Cset,\Iset$};
    \node[wide, below=of verify] (fused) {Fused forward\\$\pi_\theta^{\mathrm{mix}}=\mathrm{softmax}((1-\lambda)z^{(1)}+\lambda z^{(2)})$};
    \node[wide, below=of fused] (loss) {Stage-2 loss\\$\Ltotal^{(2)}=\Lpos^{(2)}+\beta_2\Lneg^{(2)}$};
    \node[wide, below=of loss] (update) {Joint update\\$\nabla_{\theta_1,\theta_2}\Ltotal^{(2)}$\\evaluate Model~2};

    \draw[arrow] (base) -- (stage1);
    \draw[arrow] (stage1) -- (model2);
    \draw[arrow] (base) -- (m1);
    \draw[arrow] (model2) -- (m2);
    \draw[arrow] (m2) -- (rollout);
    \draw[arrow] (rollout) -- (verify);
    \draw[arrow] (m1) |- (fused);
    \draw[arrow] (m2) |- (fused);
    \draw[arrow] (verify) -- (fused);
    \draw[arrow] (fused) -- (loss);
    \draw[arrow] (loss) -- (update);
    \draw[arrow] (update.east) .. controls +(20mm,0mm) and +(18mm,-16mm) .. (m2.south);
    \draw[arrow] (update.west) .. controls +(-20mm,0mm) and +(-18mm,-16mm) .. (m1.south);
\end{tikzpicture}
\caption{Relationship between Stage 1 and Stage 2. Stage 1 selects a single-model checkpoint that initializes Model~2. Stage 2 samples from Model~2 only, but constructs the training loss using a fused forward pass through Model~1 and Model~2, so both submodels receive gradients.}
\label{fig:stage1-stage2-flow}
\end{figure}

\section{Stage-2 Fast-Validation Results}
\label{sec:stage2-results}

\begin{table}[t]
\centering
\small
\setlength{\tabcolsep}{4.2pt}
\renewcommand{\arraystretch}{1.15}
\caption{Stage-2 fast-validation results. Both runs use Model2-policy rollout and fused-policy loss. Values are online validation percentages. ``Best checkpoint step'' is selected by MATH-500 mean@3; rows labeled ``Best'' report each metric's own peak over validation steps; ``Final'' is the step-75 validation.}
\label{tab:stage2-fast-validation}
\resizebox{\textwidth}{!}{%
\begin{tabular}{lcc}
\toprule
Result & S1 $\beta=0.0 \rightarrow$ S2 $\beta_2=0.0$ & S1 $\beta=0.1 \rightarrow$ S2 $\beta_2=0.1$ \\
\midrule
Stage-1 source checkpoint
    & step 85 & step 150 \\
Stage-1 best MATH-500 mean@3
    & 73.25 & 75.74 \\
Stage-2 latest step
    & 75 & 75 \\
Stage-2 best checkpoint step
    & 35 & 20 \\
Stage-2 best MATH-500 mean@3
    & 74.80 & \textbf{76.68} \\
Stage-2 best MATH-500 best@3
    & 80.91 & \textbf{82.24} \\
Stage-2 best AIME-2025 mean@3
    & 16.67 & \textbf{19.23} \\
Stage-2 best AIME-2025 best@3
    & 21.83 & \textbf{25.85} \\
Final MATH-500 mean@3
    & 1.41 & 1.34 \\
Final AIME-2025 mean@3
    & 0.00 & 0.00 \\
Final extraction-failure rate
    & 55.75 & 93.23 \\
Health note
    & early gain, final collapse & early gain, final collapse \\
\bottomrule
\end{tabular}
}
\end{table}

\paragraph{Interpretation.}
The Stage-2 peak metrics are consistent with the intended gradient-amplification
effect: both matched-beta runs improve over their Stage-1 source checkpoints at
early validation steps. However, both 75-step runs collapse by the final
validation, with very high answer-extraction failure rates. We therefore treat
these Stage-2 runs as useful fast-validation evidence for the method mechanism,
but not yet as stable final models. The next experimental question is how to
keep the fused-loss amplification while preventing late-format drift.

\section{Implementation Map}
\label{sec:opsft-implementation}

\begin{itemize}[leftmargin=1.5em]
    \item Loss registration: \texttt{verl/trainer/ppo/core\_algos.py}, \texttt{loss\_mode=wdl\_sft}.
    \item Reward-label override: \texttt{verl/trainer/ppo/ray\_trainer.py} writes raw verifier labels into \texttt{advantages} for \texttt{wdl\_sft}.
    \item Stage-1 launcher: \texttt{recipe/on\_policy\_wdl\_sft/staged\_v1/run\_s1\_base\_sft.sh}.
    \item Beta wrappers: \texttt{recipe/on\_policy\_wdl\_sft/staged\_v1/run\_s1\_beta\_*.sh}.
    \item Queue: \texttt{recipe/on\_policy\_wdl\_sft/staged\_v1/run\_stage1\_beta\_search\_queue.sh}.
    \item Stage-2 non-overlap shard builder: \texttt{recipe/on\_policy\_wdl\_sft/staged\_v1/create\_stage2\_nonoverlap\_shard.py}.
    \item Stage-2 fast-validation launchers: \texttt{recipe/on\_policy\_wdl\_sft/staged\_v1/run\_s2\_from\_s1\_beta*\_beta*.sh}.
    \item Stage-2 queue and monitor: \texttt{recipe/on\_policy\_wdl\_sft/staged\_v1/run\_stage2\_fast\_validation\_queue.sh} and \texttt{monitor\_stage2\_fast\_validation\_queue\_notify.sh}.
\end{itemize}

\section{Summary}
\label{sec:opsft-summary}

The current experiment is best understood as a staged progression from
verifier-selected on-policy SFT to Model2-rollout fused-loss WDL-SFT. Stage 1
uses Eq.~\eqref{eq:opsft-total-loss}: reinforce correct sampled responses,
optionally penalize incorrect sampled responses, and skip prompts with no
positive sample. Stage 2 uses Eq.~\eqref{eq:stage2-total-loss}: keep rollouts on
the current Model2 policy, but construct the supervised loss on the fused
logit distribution so gradients flow through both submodels. The first
Stage-2 results support the existence of an early gradient-amplification
effect, but the final-step validation collapse shows that the current
implementation is not yet a stable final training recipe.
